\subsection{Results}

The top panel of Figure~\ref{fig:ll_ratios} shows the estimation results.
Each panel provides the histogram of the game-by-game log-likelihood ratios ($\hat \Lambda_s$) for each comparison.
The vertical dashed black lines indicate the means of the log-likelihood ratios ($\bar \Lambda_S$), which are also given in the upper left corner of each panel.
Bootstrapped standard errors are in parentheses.
Green indicates ratios greater than zero, where the optimized AI has more predictive power, and red is the converse.

\begin{figure}[h!]
\begin{center}
\begin{minipage}{\textwidth}
\caption{Predictive power of strategic \aisub{s} compared to other models ($\varepsilon = 0.2$)} 
\label{fig:ll_ratios}
\vspace*{-0.15in}
\includegraphics[width=\textwidth]{plots/prop_ll_plots.pdf} 
\end{minipage}
\end{center}
\begin{footnotesize}
\begin{singlespace}
\vspace*{-0.05in}
\emph{Notes:} The top panel shows the distribution of the log-likelihood ratios for each benchmark.
The vertical black line indicates the mean, which is also provided in the upper left corner of each panel.
The respective bootstrapped standard errors are in parentheses.
Green indicates ratios greater than zero, where the optimized strategic AI has more predictive power, and red is the converse.
The bottom panel shows the proportion of games for which the sample of strategic \aisub{s} is the best predictor of initial play. 
Error bars are 95\% Clopper-Pearson confidence intervals
\end{singlespace}
\end{footnotesize}
\end{figure}
With $\varepsilon = 0.2$, the strategic sample of \aisub{} is, on average, significantly more predictive than any of the benchmarks ($p < 0.001$ for all comparisons).
These differences are substantial.
Starting with the leftmost panel, across all human observations in the dataset, the strategic sample of \aisub{s} achieves an average per-observation likelihood ratio of $e^{\BaselineBarLambda} = \ExpBaselineBarLambda$ in favor of the model, relative to the baseline~AI. 
In other words, the likelihood of the observed human data under the strategic \aisub{} model is, on average, $\ExpBaselineBarLambda$ times larger per observation than under the baseline (i.e., \PctBaselineBarLambda\% more probability).

Moving to the right, the corresponding average likelihood ratios are $e^{\CHBarLambda} = \ExpCHBarLambda$ compared to the cognitive hierarchy model (\PctCHBarLambda\% more probability) and $e^{\NashBarLambda} = \ExpNashBarLambda$ compared to the Harsanyi-Selten-selected equilibria ( \PctNashBarLambda\% more probability).
To be clear, this means that the strategic \aisub{s} outperform the cognitive hierarchy model by a larger margin than they outperform the Harsanyi-Selten Nash equilibria.
This is notable because throughout the literature, the cognitive hierarchy model has been a strong predictor of initial play in strategic games.
It has been explicitly shown to ``explain why equilibrium theory predicts behavior well in some games and poorly in others'' \citep{Camerer2004Cognitive}.
Of course, the cognitive hierarchy model might be more accurate with different values for $\tau$, but we had no way of knowing what these were ex ante---besides those from the literature---without the data sampled from the population we wanted to predict.
Finally, the strategic \aisub{s} show a large gain compared to the random pure-strategy benchmark $e^{\PureRandBarLambda} = \ExpPureRandBarLambda$ (\PctPureRandBarLambda\% more probability), and a more modest but still significant advantage over the uniform distribution $e^{\UniformBarLambda} = \ExpUniformBarLambda$ (\PctUniformBarLambda\% more probability).

The bottom panel of Figure~\ref{fig:ll_ratios} shows the proportion of games for which the sample of strategic \aisub{s} is the best predictor---i.e., has a positive log-likelihood ratio ($\sum_{s \in S} \mathbf 1\{\hat \Lambda_s > 0\} / |S|$).
The results are consistent with the top panel.
The theory-grounded sample of strategic \aisub{s} better predicts the human responses in more games than any other model.
This proportion is large and significant for the baseline ($\BaselineHatP$).
It is smaller, although still substantial, for both the cognitive hierarchy model ($\CHHatP$) and the Harsanyi-Selten Nash equilibria ($\NashHatP$).
Because each game has only a small number of human observations ($m_s$ between 1 and 5), sampling noise in $\hat\Lambda_s$ mechanically attenuates this proportion toward 0.5, almost surely making these estimates highly conservative.
Empirical Bayes shrinkage methods \citep[see][]{walters2024empirical} could partially correct for this, but we leave the unadjusted proportions for simplicity.
This issue does not affect the sample averages $\bar\Lambda_S$ reported above, which remain unbiased.

Tables in Appendix~\ref{app:robustness_checks} provide the same statistical analyses for $\varepsilon \in \{.05, 0.1, 0.3\}$, respectively.
The above results are robust to these additional sensitivity checks.
$\bar \Lambda_S > 0 $ for all comparisons, and the proportions of games for which the optimized \aisub{} is the best predictor are all greater than 50\%.

Beyond relative comparisons, the strategic \aisub{s} demonstrate impressive absolute predictive accuracy.
Without any smoothing, \OptHighestMatchRate\% of human respondents selected the strategy for which the optimized \aisub{} assigns the most density.
Given that the number of possible strategies per game varied evenly between 5 and 20, this is notable.
Furthermore, \OptTopThreeMatchRate\% of human respondents selected one of the top three strategies for which the strategic \aisub{} assigns the most density.
And maybe most surprisingly, for \OptGameAllMatchRate\% of games, all human respondents selected a strategy within the support of the strategic \aisub{}.
